\documentclass[11pt]{article}

\usepackage[utf8]{inputenc}
\usepackage[T1]{fontenc}
\usepackage{amsmath,amssymb,amsthm}
\usepackage{booktabs}
\usepackage{geometry}
\geometry{margin=1in}

\newtheorem{theorem}{Theorem}[section]
\newtheorem{proposition}[theorem]{Proposition}
\newtheorem{definition}[theorem]{Definition}
\newtheorem{remark}[theorem]{Remark}

\title{The Pandrosion Zeta Function\\
and the Discriminant as Spectral Determinant}
\author{Ivan Besevic}
\date{April 2026}

\begin{document}
\maketitle

\begin{abstract}
We introduce the \emph{Pandrosion zeta function}
$\zeta_P(s) := \sum_{k=1}^{d} P'(\alpha_k)^{-s}$,
an analogue of the spectral zeta function where the ``eigenvalues'' are the
derivative values $P'(\alpha_k)$ at the roots. We prove:
(i) $\zeta_P(0) = d$;
(ii) $\zeta_P(1) = 0$ for $d \geq 2$ (the \emph{vanishing identity}, proved via partial
fractions);
(iii) a hierarchy of $d-1$ vanishing relations $\sum \alpha_k^m / P'(\alpha_k) = 0$
for $m \leq d-2$;
and (iv) the discriminant is the spectral determinant:
$\exp(-\zeta'_P(0)) = \prod P'(\alpha_k) = \pm D(P)^2$.
This gives a spectral-analytic formulation of all results from Papers 20--22.
\end{abstract}

\section{Definition and basic properties}

\begin{definition}[Pandrosion zeta function]
For $P(z) = \prod_{k=1}^{d}(z-\alpha_k)$ with simple roots:
\begin{equation}
\zeta_P(s) := \sum_{k=1}^{d} P'(\alpha_k)^{-s},
\end{equation}
where $P'(\alpha_k) = Q(\alpha_k,\alpha_k) = \prod_{j \neq k}(\alpha_k - \alpha_j)$
is the Pandrosion self-value (Paper 22).
\end{definition}

\begin{proposition}[Trivial value]
$\zeta_P(0) = d$.
\end{proposition}

\section{The vanishing identity}

\begin{theorem}[Pandrosion vanishing]
For $d \geq 2$:
\begin{equation}
\boxed{\zeta_P(1) = \sum_{k=1}^{d} \frac{1}{P'(\alpha_k)} = 0.}
\end{equation}
\end{theorem}

\begin{proof}
The partial fraction decomposition of $1/P(z)$ is:
\begin{equation}
\frac{1}{P(z)} = \sum_{k=1}^{d} \frac{1}{P'(\alpha_k)\,(z - \alpha_k)}.
\end{equation}
As $z \to \infty$:
\[
\frac{1}{P(z)} = z^{-d}(1 + O(z^{-1})) \qquad \text{(no $z^{-1}$ term for $d \geq 2$).}
\]
The partial fraction sum expands as:
\[
\sum_k \frac{1}{P'(\alpha_k)(z - \alpha_k)}
= \frac{1}{z}\sum_k \frac{1}{P'(\alpha_k)} + O(z^{-2}).
\]
Matching the $z^{-1}$ coefficient: $\sum 1/P'(\alpha_k) = 0$.
\end{proof}

\begin{theorem}[Higher vanishing relations]
\begin{equation}
\sum_{k=1}^{d} \frac{\alpha_k^m}{P'(\alpha_k)} =
\begin{cases}
0 & \text{if } 0 \leq m \leq d-2,\\
1 & \text{if } m = d-1.
\end{cases}
\end{equation}
\end{theorem}

\begin{proof}
The coefficient of $z^{-1}$ in
$z^m / P(z) = z^{m-d}(1 + c_1 z^{-1} + \cdots)$
is $0$ for $m \leq d-2$ and $1$ for $m = d-1$. By partial fractions:
\[
\frac{z^m}{P(z)} = \sum_k \frac{\alpha_k^m}{P'(\alpha_k)\,(z - \alpha_k)} + \text{polynomial in } z.
\]
The coefficient of $z^{-1}$ on the right is $\sum \alpha_k^m / P'(\alpha_k)$.
\end{proof}

\begin{remark}
The $d-1$ vanishing relations form a \emph{biorthogonality system}: the vectors
$(\alpha_1^m, \ldots, \alpha_d^m)$ for $m = 0, \ldots, d-2$ are orthogonal to the
vector $(1/P'(\alpha_1), \ldots, 1/P'(\alpha_d))$. The normalisation at $m = d-1$
fixes the scale.
\end{remark}

\section{The discriminant as spectral determinant}

\begin{theorem}[Spectral determinant]
\begin{equation}
\boxed{\exp(-\zeta'_P(0)) = \prod_{k=1}^{d} P'(\alpha_k) = (-1)^{d(d-1)/2}\, D(P)^2,}
\end{equation}
where $D(P) = \prod_{i<j}(\alpha_i - \alpha_j)$ is the discriminant.
\end{theorem}

\begin{proof}
$\zeta'_P(s) = -\sum \log P'(\alpha_k) \cdot P'(\alpha_k)^{-s}$,
so $\zeta'_P(0) = -\sum \log P'(\alpha_k) = -\log \prod P'(\alpha_k)$.
Exponentiating: $\exp(-\zeta'_P(0)) = \prod P'(\alpha_k)$.

By Paper 22: $P'(\alpha_k) = \prod_{j \neq k}(\alpha_k - \alpha_j)$,
so $\prod_k P'(\alpha_k) = \prod_k \prod_{j \neq k}(\alpha_k - \alpha_j)
= (-1)^{d(d-1)/2} D(P)^2$.
\end{proof}

This is the polynomial analogue of the Minakshisundaram--Pleijel zeta function,
where the spectral determinant of a Laplacian on a manifold is computed via
$\exp(-\zeta'(0))$. Here, the ``manifold'' is the root configuration of $P$,
and the ``Laplacian eigenvalues'' are $P'(\alpha_k)$.

\section{The energy zeta function}

\begin{definition}
The \textbf{energy zeta function} is:
\begin{equation}
Z_P(s) := \sum_{k=1}^{d} E_P(\alpha_k)^{-s}
       = \sum_{k=1}^{d} P'(\alpha_k)^{-2s} = \zeta_P(2s).
\end{equation}
\end{definition}

\begin{center}
\begin{tabular}{ccl}
\toprule
$s$ & $Z_P(s)$ & Interpretation \\
\midrule
$-1$    & $\sum E_P(\alpha_k)$    & Total self-energy \\
$-1/2$  & $\sum |P'(\alpha_k)|$   & Total field strength \\
$0$     & $d$                     & Degree \\
$1/2$   & $\sum 1/|P'(\alpha_k)|$ & Sum of residue norms \\
$1$     & $\sum 1/E_P(\alpha_k)$  & ``Capacitance'' \\
\bottomrule
\end{tabular}
\end{center}

\section{Analogy with the Riemann zeta function}

\begin{center}
\begin{tabular}{lcc}
\toprule
\textbf{Property} & \textbf{Riemann $\zeta$} & \textbf{Pandrosion $\zeta_P$} \\
\midrule
Definition           & $\sum n^{-s}$                  & $\sum P'(\alpha_k)^{-s}$ \\
Value at $0$         & $-1/2$                         & $d$ \\
Vanishing?           & $\zeta(-2n)=0$ (trivial zeros) & $\zeta_P(1) = 0$ \\
Derivative at $0$    & $-\log\sqrt{2\pi}$             & $-\log |D(P)|^2$ \\
Spectral determinant & $\sqrt{2\pi}$                  & $|D(P)|^2$ \\
Euler product        & $\prod (1-p^{-s})^{-1}$        & $\prod (1 - r_k)$ (Paper 25) \\
\bottomrule
\end{tabular}
\end{center}

The discriminant $D(P)$ is the Pandrosion analogue of $\sqrt{2\pi}$: the
\emph{regularised volume} of the root configuration.

\section{Conclusion}

The Pandrosion zeta function $\zeta_P(s) = \sum P'(\alpha_k)^{-s}$ encodes
the entire spectral information of the root system in a single meromorphic
function of $s$. Its key features are:
\begin{enumerate}
\item \textbf{Vanishing at $s = 1$:} $\sum 1/P'(\alpha_k) = 0$,
a fundamental constraint on derivative values at roots.
\item \textbf{Higher vanishing:} $d - 1$ biorthogonality relations from
partial fractions.
\item \textbf{Spectral determinant:} $\exp(-\zeta'_P(0)) = \pm D(P)^2$,
giving the discriminant a zeta-theoretic meaning.
\end{enumerate}

Together with Paper 22 ($D^2 = \prod E_P(\alpha_k)$) and Paper 26
(electrostatic interpretation), this completes the analytic picture of the
Pandrosion root spectrum.

\begin{thebibliography}{9}
\bibitem{MP1949} S.~Minakshisundaram, \AA.~Pleijel,
\emph{Some properties of the eigenfunctions of the Laplace-operator},
Canad. J. Math. \textbf{1} (1949), 242--256.
\bibitem{Besevic2026} I.~Besevic,
\emph{The Pandrosion discriminant}, Preprint (2026).
\end{thebibliography}

\end{document}
